\ifdefined\gkver\else\def\gkver{student}\fi
\documentclass[\gkver]{gaokaozhenti}
\begin{document}

\chapter{2009年全国理综Ⅱ}
%% paperType: 理综物理部分

\begin{enumerate}
\item
%% number: 14
%% paperName: 2009年全国理综Ⅱ第 14 题 6 分
%% typeId: 2
%% type: 多选题
%% chapter: 机械波
%% point: 波长频率波速
%% method: 无
%% score: 6
%% degree: 650
%% duplicateId: 0
%% body:
下列关于简谐振动和简谐波的说法，正确的是\xzanswer{AD}

\fourchoices[answer=AD]
{媒质中质点振动的周期一定和相应的波的周期相等}
{媒质中质点振动的速度一定和相应的波的波速相等}
{波的传播方向一定和媒质中质点振动的方向一致}
{横波的波峰与波谷在振动方向上的距离一定是质点振幅的两倍}

%% answer: AD
%% memo:
\memoanswer{
本题考查机械波和机械振动。介质中的质点的振动周期和相应的波传播周期一致，A 正确。而各质点做简谐运动速度随时间作周期性的变化，但波在介质中是匀速向前传播的，所以不相等，B 错。对于横波而言传播方向和振动方向是垂直的，C 错。根据波的特点 D 正确。
}

\item
%% number: 15
%% paperName: 2009年全国理综Ⅱ第 15 题 6 分
%% typeId: 1
%% type: 单选题
%% chapter: 直线运动
%% point: 运动图象
%% method: 无
%% score: 6
%% degree: 450
%% duplicateId: 0
%% body:
两物体甲和乙在同一直线上运动，它们在 $0\sim0.4\Us$ 时间内的 $v-t$ 图象如图所示。若仅在两物体之间存在相互作用，则物体甲与乙的质量之比和图中时间 $t_{1}$ 分别为\xzanswer{B}

\onepicture[w=6.0cm]{figs/15.png}

\fourchoices[answer=B]
{$\frac{1}{3}$ 和 $0.30\Us$}
{$3$ 和 $0.30\Us$}
{$\frac{1}{3}$ 和 $0.28\Us$}
{$3$ 和 $0.28\Us$}

%% answer: B
%% memo:
\memoanswer{
本题考查图象问题。根据速度图象的特点可知甲做匀加速，乙做匀减速。根据 $a=\frac{\Delta v}{\Delta t}$ 得 $3a_{\text{甲}}=a_{\text{乙}}$，根据牛顿第二定律有 $\frac{F}{m_{\text{甲}}}=\frac{1}{3}\frac{F}{m_{\text{乙}}}$，得 $\frac{m_{\text{甲}}}{m_{\text{乙}}}=3$，由 $a_{\text{乙}}=\frac{4}{0.4}=10\Umsq=\frac{1}{0.4-t}$，得 $t=0.3\Us$，B 正确。
}

\item
%% number: 16
%% paperName: 2009年全国理综Ⅱ第 16 题 6 分
%% typeId: 2
%% type: 多选题
%% chapter: 气体性质
%% point: 热力学第一定律
%% method: 无
%% score: 6
%% degree: 450
%% duplicateId: 0
%% body:
如图，水平放置的密封气缸内被一竖直隔板分隔为左右两部分，隔板可在气缸内无摩擦滑动，右侧气体内有一电热丝。气缸壁和隔板均绝热。初始时隔板静止，左右两边气体温度相等。现给电热丝提供一微弱电流，通电一段时间后切断电源。当缸内气体再次达到平衡时，与初始状态相比\xzanswer{BC}

\onepicture[w=7.0cm]{figs/16.png}

\fourchoices[answer=BC]
{右边气体温度升高，左边气体温度不变}
{左右两边气体温度都升高}
{左边气体压强增大}
{右边气体内能的增加量等于电热丝放出的热量}

%% answer: BC
%% memo:
\memoanswer{
本题考查气体。当电热丝通电后，右边的气体温度升高气体膨胀，将隔板向左推，对左边的气体做功，根据热力学第一定律，内能增加，气体的温度升高。根据气体定律左边的气体压强增大。BC 正确，右边气体内能的增加值为电热丝发出的热量减去对左边的气体所做的功，D 错。
}

\item
%% number: 17
%% paperName: 2009年全国理综Ⅱ第 17 题 6 分
%% typeId: 1
%% type: 单选题
%% chapter: 电路
%% point: 电路中的图象
%% method: 无
%% score: 6
%% degree: 450
%% duplicateId: 0
%% body:
图为测量某电源电动势和内阻时得到的 $U-I$ 图线。用此电源与三个阻值均为 $3\UO$ 的电阻连接成电路，测得路端电压为 $4.8\UV$。则该电路可能为\xzanswer{B}

\onepicture[w=5.5cm]{figs/17a.png}

\fourchoices[answer=B, ispicture=true, h=2.2cm]
{figs/17b.png}
{figs/17c.png}
{figs/17d.png}
{figs/17e.png}

%% answer: B
%% memo:
\memoanswer{
本题考查测电源的电动势和内阻的实验。由测量某电源电动势和内阻时得到的 $U-I$ 图线可知该电源的电动势为 $6\UV$，内阻为 $0.5\UO$。此电源与三个均为 $3\UO$ 的电阻连接成电路时测的路端电压为 $4.8\UV$，A 中的路端电压为 $4\UV$，B 中的路端电压约为 $4.8\UV$。C 中的路端电压约为 $5.7\UV$，D 中的路端电压为 $5.4\UV$。正确选项：B。
}

\item
%% number: 18
%% paperName: 2009年全国理综Ⅱ第 18 题 6 分
%% typeId: 2
%% type: 多选题
%% chapter: 原子和原子核
%% point: 玻尔模型
%% method: 无
%% score: 6
%% degree: 450
%% duplicateId: 0
%% body:
氢原子的部分能级如图所示，已知可见光子能量在 $1.62\UeV$ 到 $3.11\UeV$ 之间。由此可推知，氢原子\xzanswer{AD}

\onepicture[w=4.0cm]{figs/18.png}

\fourchoices[answer=AD]
{从高能级向 $n=1$ 能级跃迁时发出的光的波长比可见光的短}
{从高能级向 $n=2$ 能级跃迁时发出的光为可见光}
{从高能级向 $n=3$ 能级跃迁时发出的光的频率比可见光的高}
{从 $n=3$ 能级向 $n=2$ 能级跃迁时发出的光为可见光}

%% answer: AD
%% memo:
\memoanswer{
本题考查玻尔的原子理论。从高能级向 $n=1$ 的能级跃迁的过程中辐射出的最小光子能量为 $9.20\UeV$，不在 $1.62\UeV$ 到 $3.11\UeV$ 之间，A 正确。已知可见光子能量在 $1.62\UeV$ 到 $3.11\UeV$ 之间，从高能级向 $n=2$ 能级跃迁时发出的光的能量 $\leq 3.40\UeV$，B 错。从高能级向 $n=3$ 能级跃迁时发出的光，只有能量大于 $3.11\UeV$ 的光的频率才比可见光高，C 错。从 $n=3$ 到 $n=2$ 的过程中释放的光的能量等于 $1.89\UeV$，介于 $1.62\UeV$ 到 $3.11\UeV$ 之间，所以是可见光，D 对。
}

\item
%% number: 19
%% paperName: 2009年全国理综Ⅱ第 19 题 6 分
%% typeId: 2
%% type: 多选题
%% chapter: 电场
%% point: 带电粒子的轨迹类问题
%% method: 无
%% score: 6
%% degree: 450
%% duplicateId: 0
%% body:
图中虚线为匀强电场中与场强方向垂直的等间距平行直线，两粒子 M、N 质量相等，所带电荷的绝对值也相等，现将 M、N 从虚线上的 O 点以相同速率射出，两粒子在电场中运动的轨迹分别如图中两条实线所示。点 a、b、c 为实线与虚线的交点，已知 O 点电势高于 c 点。若不计重力，则\xzanswer{BD}

\onepicture[w=6.0cm]{figs/19.png}

\fourchoices[answer=BD]
{M 带负电荷，N 带正电荷}
{N 在 a 点的速度与 M 在 c 点的速度大小相同}
{N 在从 O 点运动至 a 点的过程中克服电场力做功}
{M 在从 O 点运动至 b 点的过程中，电场力对它做的功等于零}

%% answer: BD
%% memo:
\memoanswer{
本题考查带电粒子在电场中的运动。图中的虚线为等势线，所以 M 点从 O 点到 b 点的过程中电场力对粒子做功等于零，D 正确。根据 M、N 粒子的运动轨迹可知 N 受到的电场力向上，M 受到的电场力向下，电荷的正负不清楚但为异种电荷，A 错。O 到 a 的电势差等于 O 到 c 的两点的电势差，而且电荷和质量大小相等，电场力都做的是正功，根据动能定理得 a 与 c 两点的速度大小相同，但方向不同，B 对。
}

\item
%% number: 20
%% paperName: 2009年全国理综Ⅱ第 20 题 6 分
%% typeId: 1
%% type: 单选题
%% chapter: 运动和力
%% point: 牛顿定律的应用
%% method: 无
%% score: 6
%% degree: 450
%% duplicateId: 0
%% body:
以初速度 $v_{0}$ 竖直向上抛出一质量为 $m$ 的小物块。假定物块所受的空气阻力 $f$ 大小不变。已知重力加速度为 $g$，则物体上升的最大高度和返回到原抛出点的速率分别为\xzanswer{A}

\fourchoices[answer=A]
{$\frac{v_{0}^{2}}{2g(1+f/mg)}$ 和 $v_{0}\sqrt{\frac{mg-f}{mg+f}}$}
{$\frac{v_{0}^{2}}{2g(1+f/mg)}$ 和 $v_{0}\sqrt{\frac{mg}{mg+f}}$}
{$\frac{v_{0}^{2}}{2g(1+2f/mg)}$ 和 $v_{0}\sqrt{\frac{mg-f}{mg+f}}$}
{$\frac{v_{0}^{2}}{2g(1+2f/mg)}$ 和 $v_{0}\sqrt{\frac{mg}{mg+f}}$}

%% answer: A
%% memo:
\memoanswer{
本题考查动能定理。上升的过程中，重力做负功，阻力 $f$ 做负功，由动能定理得 $-(mgh+fh)=-\frac{1}{2}mv_{0}^{2}$，$h=\frac{v_{0}^{2}}{2g(1+\frac{f}{mg})}$。求返回抛出点的速度，由全程使用动能定理重力做功为零，只有阻力做功，为 $-2fh=\frac{1}{2}mv^{2}-\frac{1}{2}mv_{0}^{2}$，解得 $v=v_{0}\sqrt{\frac{mg-f}{mg+f}}$，A 正确。
}

\item
%% number: 21
%% paperName: 2009年全国理综Ⅱ第 21 题 6 分
%% typeId: 2
%% type: 多选题
%% chapter: 光学
%% point: 全反射
%% method: 无
%% score: 6
%% degree: 450
%% duplicateId: 0
%% body:
一玻璃砖横截面如图所示，其中 ABC 为直角三角形（AC 边未画出），AB 为直角边，$\angle ABC=45^\circ$，ADC 为一圆弧，其圆心在 BC 边的中点。此玻璃的折射率为 1.5。P 为一贴近玻璃砖放置的与 AB 垂直的光屏。若一束宽度与 AB 边长度相等的平行光从 AB 边垂直射入玻璃，则\xzanswer{BD}

\onepicture[w=5.5cm]{figs/21.png}

\fourchoices[answer=BD]
{从 BC 边折射出束宽度与 BC 边长度相等的平行光}
{屏上有一亮区，其宽度小于 AB 边的长度}
{屏上有一亮区，其宽度等于 AC 边的长度}
{当屏向远离玻璃砖的方向平行移动时，屏上亮区先逐渐变小然后逐渐变大}

%% answer: BD
%% memo:
\memoanswer{
本题考查光的折射和全反射。宽为 AB 的平行光进入到玻璃中直接射到 BC 面，入射角为 $45^\circ$，大于临界角 $\theta=\arcsin\frac{1}{1.5}$，所以在 BC 面上发生全反射，仍然以宽度大小为 AB 长度的竖直向下的平行光射到 AC 圆弧面上。根据几何关系可得到在屏上的亮区宽度小于 AB 的长度，B 对，D 正确。
}

\item
%% number: 22
%% paperName: 2009年全国理综Ⅱ第 22 题 5 分
%% typeId: 4
%% type: 实验题
%% chapter: 电路
%% point: 多用表的使用
%% method: 无
%% score: 5
%% degree: 650
%% duplicateId: 0
%% body:
某同学利用多用电表测量二极管的反向电阻。完成下列测量步骤：
\begin{enumerate}
	\item
	检查多用电表的机械零点。
	\item
	将红、黑表笔分别插入正、负表笔插孔，将选择开关拨至电阻测量挡适当的量程处。
	\item
	将红、黑表笔 \tkanswer[1.4cm]{}，进行欧姆调零。
	\item
	测反向电阻时，将 \tkanswer[1.4cm]{} 表笔接二极管正极，将 \tkanswer[1.4cm]{} 表笔接二极管负极，读出电表示数。
	\item
	为了得到准确的测量结果，应让电表指针尽量指向表盘 \tkanswer[1.4cm]{}（填“左侧”、“右侧”或“中央”）；否则，在可能的条件下，应重新选择量程，并重复步骤（3）（4）。
	\item
	测量完成后，将选择开关拨向 \tkanswer[1.6cm]{} 位置。
\end{enumerate}

%% answer:
\jdanswer{
\begin{enumerate}
	\item
	（无需填写）
	\item
	（无需填写）
	\item
	短接
	\item
	红，黑
	\item
	中央
	\item
	OFF
\end{enumerate}
}

%% memo:
\memoanswer{
本题考查多用电表的使用。首先要机械调零，在选择量程后还要进行欧姆调零，而且每一次换量程都要重复这样的过程。（3）将红、黑表笔短接，即为欧姆调零。测量二极管的反向电阻时应将红笔接二极管的正极，黑笔接负极。欧姆表盘的刻度线分布不均匀，在中央的刻度线比较均匀，所以尽量让指针指向表盘的中央。测量完成后应将开关打到 OFF 挡。
}

\item
%% number: 23
%% paperName: 2009年全国理综Ⅱ第 23 题 13 分
%% typeId: 4
%% type: 实验题
%% chapter: 曲线运动
%% point: 研究平抛运动实验
%% method: 无
%% score: 13
%% degree: 450
%% duplicateId: 0
%% body:
某同学利用图 \ref{2009全国理综Ⅱ23a} 所示装置做“研究平抛运动”的实验，根据实验结果在坐标纸上描出了小球水平抛出后的运动轨迹，但不慎将画有轨迹图线的坐标纸丢失了一部分，剩余部分如图 \ref{2009全国理综Ⅱ23b} 所示。图 \ref{2009全国理综Ⅱ23b} 中水平方向与竖直方向每小格的长度均代表 $0.10\Um$，$P_{1}$、$P_{2}$ 和 $P_{3}$ 是轨迹图线上的 3 个点，$P_{1}$ 和 $P_{2}$、$P_{2}$ 和 $P_{3}$ 之间的水平距离相等。

完成下列填空：（重力加速度取 $9.8\Umsq$）
\begin{enumerate}
	\item
	设 $P_{1}$、$P_{2}$ 和 $P_{3}$ 的横坐标分别为 $x_{1}$、$x_{2}$ 和 $x_{3}$，纵坐标分别为 $y_{1}$、$y_{2}$ 和 $y_{3}$，从图 \ref{2009全国理综Ⅱ23b} 中可读出 $|y_{1}-y_{2}|=$ \tkanswer[1.6cm]{} $\Um$，$|y_{1}-y_{3}|=$ \tkanswer[1.6cm]{} $\Um$，$|x_{1}-x_{2}|=$ \tkanswer[1.6cm]{} $\Um$（保留两位小数）。
	\item
	若已测知抛出后小球在水平方向做匀速运动，利用（1）中读取的数据，求小球从 $P_{1}$ 运动到 $P_{2}$ 所用的时间为 \tkanswer[1.4cm]{} $\Us$，小球抛出后的水平速度为 \tkanswer[1.6cm]{} $\Ums$（均可用根号表示）。
	\item
	已测得小球抛出前下滑的高度为 $0.50\Um$，设 $E_{1}$ 和 $E_{2}$ 分别为开始下滑时和抛出时的机械能，则小球从开始下滑到抛出的过程中机械能的相对损失 $\frac{E_{1}-E_{2}}{E_{1}}\times100\%=$ \tkanswer[1.4cm]{}\%（保留两位有效数字）。
\end{enumerate}
\twopicture[numstyle=arabic, labelA=2009全国理综Ⅱ23a, labelB=2009全国理综Ⅱ23b, wA=6.0cm, wB=5.0cm, align=right]
{figs/23a.png}
{figs/23b.png}

%% answer:
\jdanswer{
\begin{enumerate}
	\item
	$0.60$，$1.60$，$0.60$
	\item
	$0.20$，$3.0$
	\item
	$8.2$
\end{enumerate}
}

%% memo:
\memoanswer{
本题考查研究平抛运动的实验。由图可知 $P_{1}$ 到 $P_{2}$ 两点在竖直方向的间隔为 6 格，$P_{1}$ 到 $P_{3}$ 两点在竖直方向的间隔为 16 格，所以有 $|y_{1}-y_{2}|=0.60\Um$，$|y_{1}-y_{3}|=1.60\Um$。$P_{1}$ 到 $P_{2}$ 两点在水平方向的距离为 6 个格，则有 $|x_{1}-x_{2}|=0.60\Um$。

（2）由水平方向的运动特点可知 $P_{1}$ 到 $P_{2}$ 与 $P_{2}$ 到 $P_{3}$ 的时间相等，根据 $\Delta x=at^{2}$，解得时间约为 $0.2\Us$，则有 $v_{0}=\frac{x}{t}=\frac{0.60}{0.2}=3.0\Ums$。

（3）设抛出点为势能零点，则开始下滑时的机械能为 $E_{1}=mgh=mg/2$，抛出时的机械能为 $E_{2}=\frac{1}{2}mv_{0}^{2}=4.5m$，则根据 $\frac{E_{1}-E_{2}}{E_{1}}=0.082$。
}

\item
%% number: 24
%% paperName: 2009年全国理综Ⅱ第 24 题 15 分
%% typeId: 6
%% type: 计算题
%% chapter: 电磁感应
%% point: 电磁感应电路分析
%% method: 无
%% score: 15
%% degree: 650
%% duplicateId: 0
%% body:
如图，匀强磁场的磁感应强度方向垂直于纸面向里，大小随时间的变化率 $\frac{\Delta B}{\Delta t}=k$，$k$ 为负的常量。用电阻率为 $\rho$、横截面积为 $S$ 的硬导线做成一边长为 $l$ 的方框，将方框固定于纸面内，其右半部位于磁场区域中，求：
\begin{enumerate}
	\item
	导线中感应电流的大小；
	\item
	磁场对方框作用力的大小随时间的变化。
\end{enumerate}
\onepicture[w=5.5cm, align=right]{figs/24.png}

%% answer:
\jdanswer{
\begin{enumerate}
	\item
	$I=\frac{klS}{8\rho}$
	\item
	$\frac{k^{2}l^{2}S}{8\rho}$
\end{enumerate}
}

%% memo:
\memoanswer{
本题考查电磁感应现象。

（1）线框中产生的感应电动势 $\varepsilon=\frac{\Delta\phi}{\Delta t}=\frac{\Delta B\cdot S'}{\Delta t}=\frac{1}{2}l^{2}k$ ①

线框中产生的感应电流 $I=\frac{\varepsilon}{R}$ ②

$R=\rho\frac{4l}{S}$ ③

联立①②③得 $I=\frac{klS}{8\rho}$。

（2）导线框所受磁场力的大小为 $F=BIl$，它随时间的变化率为 $\frac{\Delta F}{\Delta t}=Il\frac{\Delta B}{\Delta t}$，由以上各式联立可得 $\frac{\Delta F}{\Delta t}=\frac{k^{2}l^{2}S}{8\rho}$。
}

\item
%% number: 25
%% paperName: 2009年全国理综Ⅱ第 25 题 18 分
%% typeId: 6
%% type: 计算题
%% chapter: 磁场
%% point: 带电粒子在组合场中的运动
%% method: 无
%% score: 18
%% degree: 250
%% duplicateId: 0
%% body:
如图，在宽度分别为 $l_{1}$ 和 $l_{2}$ 的两个毗邻的条形区域分别有匀强磁场和匀强电场，磁场方向垂直于纸面向里，电场方向与电、磁场分界线平行。一带正电荷的粒子以速率 $v$ 从磁场区域上边界的 P 点斜射入磁场，然后以垂直于电、磁场分界线的方向进入电场，最后从电场边界上的 Q 点射出，已知 PQ 垂直于电场方向，粒子轨迹与电、磁场分界线的交点到 PQ 的距离为 $d$。不计重力，求电场强度与磁感应强度大小之比及粒子在磁场与电场中运动时间之比。

\onepicture[w=5.5cm, align=right]{figs/25.png}

%% answer:
\jdanswer{$\frac{E}{B}=\frac{l_{1}^{2}+d^{2}}{l_{2}^{2}}v$；$\frac{t_{1}}{t_{2}}=\frac{l_{1}^{2}+d^{2}}{2dl_{2}}\arcsin\left(\frac{2dl_{1}}{l_{1}^{2}+d^{2}}\right)$}

%% memo:
\memoanswer{
本题考查带电粒子在有界磁场中的运动。

粒子在磁场中做匀速圆周运动，如图所示。由于粒子在分界线处的速度与分界线垂直，圆心 O 应在分界线上，OP 长度即为粒子运动的圆弧的半径 $R$。由几何关系得
$R^{2}=l_{1}^{2}+(R-d)^{2}$ ①

设粒子的质量和所带正电荷分别为 $m$ 和 $q$，由洛伦兹力公式和牛顿第二定律得
$qvB=m\frac{v^{2}}{R}$ ②

设 $P'$ 为虚线与分界线的交点，$\angle POP'=\alpha$，则粒子在磁场中的运动时间为
$t_{1}=\frac{R\alpha}{v}$ ③

式中有 $\sin\alpha=\frac{l_{1}}{R}$ ④

粒子进入电场后做类平抛运动，其初速度为 $v$，方向垂直于电场。

设粒子的加速度大小为 $a$，由牛顿第二定律得 $qE=ma$ ⑤

由运动学公式有 $d=\frac{1}{2}at^{2}$ ⑥，$l_{2}=vt_{2}$ ⑦

由①②⑤⑥⑦式得 $\frac{E}{B}=\frac{l_{1}^{2}+d^{2}}{l_{2}^{2}}v$ ⑧

由①③④⑦式得 $\frac{t_{1}}{t_{2}}=\frac{l_{1}^{2}+d^{2}}{2dl_{2}}\arcsin\left(\frac{2dl_{1}}{l_{1}^{2}+d^{2}}\right)$

\onepicture[w=7.0cm]{figs/25_an.png}
}

\item
%% number: 26
%% paperName: 2009年全国理综Ⅱ第 26 题 21 分
%% typeId: 6
%% type: 计算题
%% chapter: 曲线运动
%% point: 万有引力定律
%% method: 对称法
%% score: 21
%% degree: 250
%% duplicateId: 0
%% body:
如图，P、Q 为某地区水平地面上的两点，在 P 点正下方一球形区域内储藏有石油，假定区域周围岩石均匀分布，密度为 $\rho$，石油密度远小于 $\rho$，可将上述球形区域视为空腔。如果没有这一空腔，则该地区重力加速度（正常值）沿竖直方向；当存在空腔时，该地区重力加速度的大小和方向会与正常情况有微小偏差，重力加速度在原竖直方向（即 PO 方向）上的投影相对于正常值的偏离叫做“重力加速度反常”。为了探寻石油区域的位置和石油储量，常利用 P 点附近重力加速度反常现象，已知引力常数为 $G$。
\begin{enumerate}
	\item
	设球形空腔体积为 $V$，球心深度为 $d$（远小于地球半径），$PQ=x$，求空腔所引起的 Q 点处的重力加速度反常；
	\item
	若在水平地面上半径 $L$ 的范围内发现：重力加速度反常值在 $\delta$ 与 $k\delta$（$k>1$）之间变化，且重力加速度反常的最大值出现在半径为 $L$ 的范围的中心，如果这种反常是由于地下存在某一球形空腔造成的，试求此球形空腔球心的深度和空腔的体积。
\end{enumerate}
\onepicture[w=5.5cm, align=right]{figs/26.png}

%% answer:
\jdanswer{
\begin{enumerate}
	\item
	$\Delta g'=\frac{G\rho Vd}{(d^{2}+x^{2})^{3/2}}$
	\item
	$d=\frac{L}{\sqrt{k^{2/3}-1}}$，$V=\frac{L^{2}k\delta}{G\rho(k^{2/3}-1)}$
\end{enumerate}
}

%% memo:
\memoanswer{
本题考查万有引力部分的知识。

（1）如果将近地表的球形空腔填满密度为 $\rho$ 的岩石，则该地区重力加速度便回到正常值。因此，重力加速度反常可通过填充后的球形区域产生的附加引力 $G\frac{Mm}{r^{2}}=m\Delta g$ ① 来计算，式中的 $m$ 是 Q 点处某质点的质量，$M$ 是填充后球形区域的质量，$M=\rho V$ ②。

而 $r$ 是球形空腔中心 O 至 Q 点的距离 $r=\sqrt{d^{2}+x^{2}}$ ③。$\Delta g$ 在数值上等于由于存在球形空腔所引起的 Q 点处重力加速度改变的大小。Q 点处重力加速度改变的方向沿 OQ 方向，重力加速度反常 $\Delta g'$ 是这一改变在竖直方向上的投影 $\Delta g'=\frac{d}{r}\Delta g$ ④。联立以上式子得
$\Delta g'=\frac{G\rho Vd}{(d^{2}+x^{2})^{3/2}}$ ⑤

（2）由⑤式得，重力加速度反常 $\Delta g'$ 的最大值和最小值分别为
$\left(\Delta g'\right)_{\max}=\frac{G\rho V}{d^{2}}$ ⑥

$\left(\Delta g'\right)_{\min}=\frac{G\rho Vd}{(d^{2}+L^{2})^{3/2}}$ ⑦

由题设有 $\left(\Delta g'\right)_{\max}=k\delta$、$\left(\Delta g'\right)_{\min}=\delta$ ⑧

联立以上式子得，地下球形空腔球心的深度和空腔的体积分别为
$d=\frac{L}{\sqrt{k^{2/3}-1}}$，$V=\frac{L^{2}k\delta}{G\rho(k^{2/3}-1)}$。
}

\end{enumerate}

\end{document}
